
% ============================================================
% Research-track block: Route A preregistered derivation attempt
% ============================================================
\subsection{Route A Line-Density Target: Preregistered Derivation and Falsification Harness}
\label{sec:routeA_preregistered_falsification}

\paragraph{Status.}
This subsection is \textbf{Research Track} and is explicitly marked
\textbf{not derived}.  It records a preregistered falsification attempt for
Route A, after the look-elsewhere audit showed that the earlier A--D formulas
are algebraically one seed relation rather than four independent derivations.

\paragraph{Route-A target.}
The horizon-piercing entropy ansatz is
\begin{equation}
  \langle N_{\rm pierce}\rangle = {1\over 2}\Lambda_L A,
  \qquad
  S_{\rm pierce}=k_B\sigma_{\rm pierce}N_{\rm pierce},
\end{equation}
so comparison with the Bekenstein--Hawking area coefficient gives the target
\begin{equation}
  \sigma_{\rm pierce}\Lambda_L
  = {1\over 2L_p^2}
  = 1.914036558579\times10^{69}\ {\rm m^{-2}}.
\end{equation}
This target is not a derivation; it is the number Route A must reproduce
without using $G$, $L_p$, or $t_p$ as inputs.

\paragraph{Parallel model families.}
Three preregistered families were tested.
First, an Onsager/KT-like vortex-gas estimate with core cutoff $a=r_c$ gives
at most $\sigma\Lambda\sim r_c^{-2}$, or $(\pi r_c^2)^{-1}$ with disk packing,
and misses the target by roughly forty orders of magnitude.  Equivalently, the
required multiplier
\begin{equation}
  y_{\rm req}=(\sigma\Lambda)_{\rm target}r_c^2
  \simeq 3.80\times10^{39}
\end{equation}
is not an ordinary BKT fugacity.

Second, Crofton/stereology derives only the projection factor
\begin{equation}
  \langle |\cos\theta|\rangle_{S^2}={1\over2},
\end{equation}
and therefore supports
\begin{equation}
  \langle N_{\rm pierce}\rangle/A={1\over2}\Lambda_L.
\end{equation}
It does not determine the scale of $\Lambda_L$.

Third, a torsion-channel phase-space count using a strict Weyl coefficient for
two transverse channels gives
\begin{equation}
  (\sigma\Lambda)_{\rm Weyl3}
  =r_c^{-2}{\rho_{\rm core}\over\rho_f}{2\over6\pi^2}
   \left({c\over\vchar}\right)^3,
\end{equation}
or in a paired-channel variant
\begin{equation}
  (\sigma\Lambda)_{\rm WeylPair}
  =r_c^{-2}{\rho_{\rm core}\over\rho_f}
   \left({2\over6\pi^2}\right)^2
   \left({c\over\vchar}\right)^6.
\end{equation}
The paired count gets the exponent six without a coefficient scan but remains
about $1.4\times10^3$ below the target.  The earlier seed coefficient
$16/\pi^2$ is therefore retained only as a fitted negative-control relation,
not as evidence.

\paragraph{Conclusion.}
No non-fitted model in this preregistered set derives the Route-A target.  The
only canonically safe result is the stereological factor $1/2$ and the sharp
open target for $\sigma_{\rm pierce}\Lambda_L$.  Canonization requires an
independent derivation of either $\Lambda_L$, $\sigma_{\rm pierce}$, or a
non-fitted torsion-channel degeneracy.
